
\subsection{Probability}

The probability inference agent is motivated by the imperfect-information
problem,and it uses a cheape one-step response model instead of explicitly sampling many complete hidden states.
The key question is:
after we play an action \(a\), how likely is the next player to beat it, pass, or
spend a strong response such as a bomb? This idea is implemented by
\texttt{probability\_inference.py} and used by
\texttt{probabilistic\_response\_agent.py}.

Given an information set \(I\), the agent first constructs the remaining unknown
card multiset \(U(I)\) by removing its own hand and all public played cards from
the full deck. If the next player has \(m\) cards left, then a possible hidden
hand is treated as a size-\(m\) draw from \(U(I)\). For a candidate response
\(r\), let \(R(r)\) be the multiset of cards required to play that response. The
probability that the responder can physically hold this response is estimated by
a multivariate hypergeometric count:

\[
    p_{\mathrm{has}}(r)
    =
    \frac{
        \#\{h:\ |h|=m,\ R(r)\subseteq h,\ h\subseteq U(I)\}
    }{
        \binom{|U(I)|}{m}
    }.
\]

In the code, this is computed by \texttt{contains\_probability}. It uses dynamic
programming over card ranks rather than sampling hidden hands, so the estimate is
deterministic and relatively fast. Candidate responses are generated by
\texttt{possible\_response\_actions}: it enumerates legal symbolic actions that
could beat the target action, including same-type higher responses, bombs, and
rocket.

The holding probability alone is not enough, because a player may choose not to
beat even when they can. Therefore, \texttt{estimate\_response\_distribution}
combines card availability with a simple behavioral model. Each candidate
response receives an unnormalized weight

\[
    w(r)
    =
    p_{\mathrm{has}}(r)
    \exp\left(\frac{s(r,I)}{\tau}\right),
\]

where \(s(r,I)\) is a heuristic attractiveness score and \(\tau\) is a
temperature parameter. The score prefers cheap minimal responses, gives a large
bonus when the response can finish the responder's hand, encourages enemies to
beat dangerous actions, and discourages teammates from beating their own side
unless they can finish. Bombs and rocket are penalized unless the situation is
urgent.

Pass is also modeled explicitly. If \(p_{\mathrm{beat}}\) is the approximate
probability that the responder has at least one beating response, then the raw
pass mass is

\[
    w(\mathrm{pass})
    =
    (1-p_{\mathrm{beat}})
    +
    p_{\mathrm{beat}}\,\rho_{\mathrm{pass}},
\]

where \(\rho_{\mathrm{pass}}\) is a role-dependent strategic pass parameter.
For example, enemies have a lower strategic pass rate, while teammates have a
higher one because they often should not block the current player's move. The
final response distribution is normalized as

\[
    P(r\mid I,a)
    =
    \frac{w(r)}
    {w(\mathrm{pass})+\sum_{r'\in\mathcal{R}(I,a)}w(r')}.
\]

The agent then evaluates each root action by combining an immediate action score
with the expected value of the next player's response:

\begin{align*}
    Q(a)
    = & \ S_{\mathrm{imm}}(a,I) \\
      & +
    \lambda
    \sum_{r\in\mathcal{R}(I,a)\cup\{\mathrm{pass}\}}
    P(r\mid I,a)\,V_{\mathrm{resp}}(a,r,I).
\end{align*}

Here \(S_{\mathrm{imm}}\) rewards finishing, reducing hand badness, playing
useful structures such as chains or trios, and releasing low cards, while
penalizing unnecessary high cards, bombs, and bad passes. The response value
\(V_{\mathrm{resp}}\) is negative when an enemy is likely to beat the action,
especially if the enemy can finish, but positive when the enemy passes and the
agent keeps initiative. For teammate responses, the model gives a large bonus if
the teammate can finish and otherwise discourages unnecessary interference.

To avoid making the policy worse than a strong simple baseline, the probability
agent uses the greedy RLCard policy as an anchor. It first obtains the greedy
action, then scores a pruned set of candidate root actions using the probability
model. The probability-based action is allowed to replace the greedy action only
in controlled cases, such as immediate finishing, clearly better leading
structures, or stopping a dangerous enemy. This makes the method a fast
probabilistic one-step lookahead policy rather than a full POMDP solver: it does
not search many future rounds, but it directly models the most important hidden
variable for the next decision, namely whether the next player can and will
respond.
